\section{补充对照与参数明细}\label{app:checks}
\subsection{设备与预测参数检验}\label{app:parameter-checks}
表\ref{tab:q1-sensitivity}列出问题一的设备扰动结果，各组保持相同输入曲线与6000\,kWh首末库存，每次只改变一个参数。
\par
\begin{table}[!htbp]
\centering
\caption{问题一的单参数扰动与最低电费}\label{tab:q1-sensitivity}
\normalsize
\begin{tabular}{lrrr}
\toprule
\multicolumn{1}{c}{扰动参数} & \multicolumn{1}{c}{取值} & \multicolumn{1}{c}{费用/元} & \multicolumn{1}{c}{费用变化/\%} \\
\midrule
可用容量/kWh & 7680 & 36,704.73 & 4.49 \\
可用容量/kWh & 9600 & 35,126.95 & 0.00 \\
可用容量/kWh & 11520 & 33,868.53 & -3.58 \\
功率上限/kW & 4000 & 35,247.87 & 0.34 \\
功率上限/kW & 5000 & 35,126.95 & 0.00 \\
功率上限/kW & 6000 & 35,083.29 & -0.12 \\
单向效率 & 0.85 & 36,603.41 & 4.20 \\
单向效率 & 0.9 & 35,126.95 & 0.00 \\
单向效率 & 0.95 & 33,767.03 & -3.87 \\
\bottomrule
\end{tabular}
\end{table}
\par

\FloatBarrier
表\ref{tab:hourly-subsets}保留全部八种预报子集；各组保持主参数、三次重规划及共同热启动，仅改变允许读取的新预报。
% Generated by scripts/build_paper_tables.py; do not edit by hand.
\par
\begin{table}[htbp]
\centering
\caption{问题3新小时光伏预报的完整子集对照}\label{tab:hourly-subsets}
\small
\setlength{\tabcolsep}{4pt}
\renewcommand{\arraystretch}{1.15}
\begin{tabular}{lrrr}
\toprule
新增预报时刻 & 纯附件3费用（元） & 融合策略费用（元） & 融合紧急电（kWh） \\
\midrule
无 & 13,580,285.75 & 13,333,061.82 & 29,178.88 \\
18 & 13,580,285.30 & 13,333,062.82 & 29,178.86 \\
12 & 13,564,801.30 & 13,344,238.23 & 29,324.32 \\
12、18 & 13,564,801.43 & 13,344,238.54 & 29,324.42 \\
06 & 13,585,837.75 & 13,335,964.19 & 25,744.14 \\
06、18 & 13,585,837.09 & 13,335,963.52 & 25,744.09 \\
06、12 & 13,511,207.99 & 13,315,835.68 & 25,862.47 \\
06、12、18 & 13,511,208.12 & 13,315,835.99 & 25,862.57 \\
\bottomrule
\end{tabular}
\par\smallskip
\begin{minipage}{0.97\linewidth}\footnotesize 各行均保留06、12、18点重规划、实际SOC、负载修正及当前光伏观测锚点；仅开关新小时预报。融合子集共用1月热启动。\end{minipage}
\end{table}
\par

表\ref{tab:q3-reserve}比较日末预测恢复目标，费用差伴随边界库存变化。
\par
\begin{table}[!htbp]
\centering
\caption{日末预测储备目标的变化（2—12月）}\label{tab:q3-reserve}
\normalsize
\begin{tabular}{lrrr}
\toprule
\multicolumn{1}{c}{目标/kWh} & \multicolumn{1}{c}{总费用/元} & \multicolumn{1}{c}{期初储电量/kWh} & \multicolumn{1}{c}{期末储电量/kWh} \\
\midrule
3000 & 13,298,180.06 & 4,231.74 & 3,105.48 \\
6000 & 13,315,835.99 & 7,089.87 & 5,987.72 \\
9000 & 13,346,622.47 & 9,942.13 & 8,916.25 \\
\bottomrule
\end{tabular}
\end{table}
\par

\FloatBarrier
表\ref{tab:q3-forecast}、\ref{tab:q4-3-forecast}分别比较固定与波动电价下的光伏来源，各组使用共同热启动；表\ref{tab:q4-calibration}为1月15—31日选参费用，问题4-3固定融合权重0.50。
\par
\begin{table}[!htbp]
\centering
\caption{问题三的光伏预测基准（共同热启动，2—12月）}\label{tab:q3-forecast}
\normalsize
\begin{tabular}{lrrr}
\toprule
\multicolumn{1}{c}{预测来源} & \multicolumn{1}{c}{$\lambda$} & \multicolumn{1}{c}{$\alpha$} & \multicolumn{1}{c}{总费用/元} \\
\midrule
纯历史 & 0 & 0.65 & 13,428,215.91 \\
纯附件三 & 1 & 0.65 & 13,511,208.12 \\
融合 & 0.5 & 0.65 & 13,315,835.99 \\
\bottomrule
\end{tabular}
\end{table}
\par

\par
\begin{table}[!htbp]
\centering
\caption{问题4-3的光伏预测基准（共同热启动，2—12月）}\label{tab:q4-3-forecast}
\normalsize
\begin{tabular}{lrrr}
\toprule
\multicolumn{1}{c}{预测来源} & \multicolumn{1}{c}{$\lambda$} & \multicolumn{1}{c}{$\alpha$} & \multicolumn{1}{c}{总费用/元} \\
\midrule
纯历史 & 0 & 0.65 & 14,166,961.16 \\
纯附件三 & 1 & 0.65 & 14,258,363.92 \\
融合 & 0.5 & 0.65 & 14,044,885.25 \\
\bottomrule
\end{tabular}
\end{table}
\par

\par
\begin{table}[!htbp]
\centering
\caption{波动电价的1月分位评分（1月15—31日）}\label{tab:q4-calibration}
\normalsize
\begin{tabular}{lrr}
\toprule
\multicolumn{1}{c}{$\alpha$} & \multicolumn{1}{c}{问题4-2费用/元} & \multicolumn{1}{c}{问题4-3费用/元} \\
\midrule
0.5 & 1,029,631.23 & 964,022.95 \\
0.65 & 985,076.81 & 957,308.74 \\
0.7 & 974,056.08 & 959,947.92 \\
0.8 & 966,879.67 & 972,861.43 \\
0.9 & 991,846.47 & 996,125.45 \\
0.95 & 1,013,304.20 & 1,015,980.29 \\
\bottomrule
\end{tabular}
\end{table}
\par

\FloatBarrier

\subsection{逐笔收费的替代合约}\label{sec:contracts}
若日内增购成交后不可免费撤回，需逐次记录修订。令$q^{(0)}=g$，第$r$次修订为$q^{(r)}$，则在适用交付段收取
\begin{equation}
 \begin{aligned}[b]
 C^{\mathrm{trade}}&=\sum_{d,t}\Biggl[p_{d,t}g_{d,t}+5p_{d,t}e_{d,t}\\
 &\qquad+\sum_r\left(1.5p_{d,t}\pos{q^{(r)}_{d,t}-q^{(r-1)}_{d,t}}
 +0.5p_{d,t}\pos{q^{(r-1)}_{d,t}-q^{(r)}_{d,t}}\right)\Biggr].
 \end{aligned}
\end{equation}
这里$p_{d,t}$仍按交付段结算。因原购电款不返还且余电可以不利用，减购先前承诺没有优势，替代规划采用$q^{(r)}_t\ge q^{(r-1)}_t$，只将新增量逐笔计费。

表\ref{tab:contracts}区分主合约、冻结主调度的逐笔重计、逐笔制度下重新优化三种情形。冻结重计只更换账单规则，保留原调度与库存；重优化则从1月重新选参和运行，会改变热启动与后续路径，参数和首末库存见表\ref{tab:contract-parameters}。
% Generated by scripts/build_paper_tables.py; do not edit by hand.
\par
\begin{table}[htbp]
\centering
\caption{合约解释与相应策略费用（2—12月，元）}\label{tab:contracts}
\small
\setlength{\tabcolsep}{4pt}
\renewcommand{\arraystretch}{1.15}
\begin{tabular}{lrrr}
\toprule
方案 & 结算与策略 & 总费用 & 较无调整方案节省 \\
\midrule
问题3 & 主合约 & 13,315,835.99 & 592,339.96 \\
问题3 & 冻结主调度，逐笔重计 & 13,642,823.29 & 265,352.66 \\
问题3 & 逐笔合约重新优化 & 13,460,502.62 & 447,673.33 \\
问题4-3 & 主合约 & 14,044,885.25 & 630,959.63 \\
问题4-3 & 冻结主调度，逐笔重计 & 14,390,498.18 & 285,346.69 \\
问题4-3 & 逐笔合约重新优化 & 14,199,964.36 & 475,880.51 \\
\bottomrule
\end{tabular}
\par\smallskip
\begin{minipage}{0.97\linewidth}\footnotesize 冻结重计保持调度路径；重新优化另以1月选参，并从1月起使用逐笔合约，因而也改变状态路径。\end{minipage}
\end{table}
\par

% Generated by scripts/build_paper_tables.py; do not edit by hand.
\par
\begin{table}[!htbp]
\centering
\caption{逐笔合约重优化的参数与库存（电量单位：kWh）}\label{tab:contract-parameters}
\normalsize
\begin{tabular}{lrrrrr}
\toprule
\multicolumn{1}{c}{方案} & \multicolumn{1}{c}{$\lambda$} & \multicolumn{1}{c}{$\alpha$} & \multicolumn{1}{c}{紧急购电} & \multicolumn{1}{c}{\shortstack{2月初\\储电量}} & \multicolumn{1}{c}{\shortstack{年末\\储电量}} \\
\midrule
问题3 & 0.25 & 0.65 & 27,362.36 & 7,089.87 & 5,987.72 \\
问题4-3 & 0.25 & 0.65 & 27,071.28 & 7,114.54 & 5,997.11 \\
\bottomrule
\end{tabular}
\end{table}
\par

\FloatBarrier
重新优化后，固定与波动价格下分别比对应无调整方案节省447673.33元、475880.51元，均低于主合约的节省。冻结重计费用更高，说明合约解释会影响经济结论；但其差额不能替代新制度下重新优化的结果。

\subsection{价格感知储能反馈对照}\label{sec:feedback}
为检验优先补足当前缺口的反馈是否足够，在每个十分钟段固定剩余常规承诺$q$，使用最新已允许时刻的预测，将首段净负荷和价格换为当段已知真值，求解
\begin{equation}
 \min\sum_{k=t}^{143}5\widehat p_k e_k,\qquad
 c_k\le\pos{q_k-\widetilde N_k},\quad b_k\le\pos{\widetilde N_k-q_k}.
\end{equation}
同时满足能量平衡、功率和库存边界，并以真实库存为初值。该模型只以常规富余电充电，紧急电只补缺口，午夜库存价值取0；每次只执行第一段动作。只有到允许的发布时间，才更新完整预测并重申报常规计划。

两种反馈共用参数、预报发布时间、1月热启动及正式期初库存，后续计划随各自真实状态更新。表\ref{tab:mpc-comparison}比较常规费、紧急费与总费，表\ref{tab:mpc-inventory}报告费用增量和边界库存；增费定义为滚动反馈总费减去原反馈总费。
% Generated by scripts/build_paper_tables.py; do not edit by hand.
\par
\begin{table}[htbp]
\centering
\caption{确定性价格感知反馈与原反馈的费用比较（2—12月，元）}\label{tab:mpc-comparison}
\footnotesize
\setlength{\tabcolsep}{4pt}
\renewcommand{\arraystretch}{1.15}
\begin{tabular}{lrrrr}
\toprule
方案 & 反馈规则 & 常规购电费 & 紧急购电费 & 总费用 \\
\midrule
问题2 & 贪心 & 13,291,625.72 & 616,550.23 & 13,908,175.95 \\
问题2 & 确定性滚动 & 13,287,773.99 & 654,732.22 & 13,942,506.21 \\
问题3 & 贪心 & 13,177,405.10 & 138,430.89 & 13,315,835.99 \\
问题3 & 确定性滚动 & 13,062,906.94 & 466,563.01 & 13,529,469.95 \\
问题4-2 & 贪心 & 14,032,328.84 & 643,516.03 & 14,675,844.87 \\
问题4-2 & 确定性滚动 & 14,028,077.93 & 681,762.73 & 14,709,840.66 \\
问题4-3 & 贪心 & 13,917,214.62 & 127,670.62 & 14,044,885.25 \\
问题4-3 & 确定性滚动 & 13,807,269.57 & 451,105.26 & 14,258,374.82 \\
\bottomrule
\end{tabular}
\par\smallskip
\begin{minipage}{0.97\linewidth}\footnotesize 常规购电费为原计划费与调整费之和。两种反馈共用参数、发布时间、1月热启动和正式期初库存；反馈改变库存后，后续常规计划也随之改变。\end{minipage}
\end{table}
\par

% Generated by scripts/build_paper_tables.py; do not edit by hand.
\par
\begin{table}[htbp]
\centering
\caption{滚动反馈的费用增量与物理统计}\label{tab:mpc-inventory}
\footnotesize
\setlength{\tabcolsep}{4pt}
\renewcommand{\arraystretch}{1.15}
\begin{tabular}{lrrrrr}
\toprule
方案 & 增费（元） & 紧急电（kWh） & 未利用电（kWh） & 期初SOC & 期末SOC \\
\midrule
问题2 & 34,330.27 & 147,630.86 & 2,692,215.80 & 6,840.33 & 6,258.37 \\
问题3 & 213,633.96 & 128,007.58 & 2,053,973.08 & 7,089.87 & 5,987.72 \\
问题4-2 & 33,995.79 & 143,395.60 & 2,674,373.02 & 6,864.27 & 6,263.02 \\
问题4-3 & 213,489.58 & 120,113.54 & 2,040,532.00 & 7,114.54 & 5,997.11 \\
\bottomrule
\end{tabular}
\par\smallskip
\begin{minipage}{0.97\linewidth}\footnotesize 首末SOC单位为kWh，且分别与对应原反馈相同。增费为滚动反馈总费减去原反馈总费。\end{minipage}
\end{table}
\par

\FloatBarrier
四组滚动反馈均降低常规费，却增加了更多紧急费；期末库存分别与原反馈一致，因而费用差不能归于不同的边界库存。本对照没有重新校准参数或设置跨日库存价值，只能说明该配置未改善累计费用，不能证明原反馈最优。

\subsection{上月选参的滚动评价}\label{sec:rolling}
在2025年4—12月275天另作时序实验：每月开始前，仅用上一个完整月的费用在30个候选中重选融合权重与分位；历史光伏基准只在$\lambda=0$下选分位。候选评分均从同一参考月初库存开始，不添加末端库存估值，也不将候选月底库存传入正式路径。

固定参数、逐月融合、逐月历史三条正式路径共用4月1日库存，之后独立连续运行。累计账单与首末库存见表\ref{tab:rolling-summary}，其4—12月费用不能直接与2—12月主结果相减。
% Generated by scripts/build_paper_tables.py; do not edit by hand.
\par
\begin{table}[htbp]
\centering
\caption{按月选参与固定参数比较（2025年4—12月）}\label{tab:rolling-summary}
\footnotesize
\setlength{\tabcolsep}{4pt}
\renewcommand{\arraystretch}{1.15}
\begin{tabular}{lrrrrr}
\toprule
方案 & 参数规则 & 总费用（元） & 紧急电（kWh） & 期初SOC & 期末SOC \\
\midrule
问题3 & 固定1月参数 & 11,025,848.81 & 24,392.62 & 6,461.68 & 5,987.72 \\
问题3 & 上月选择融合 & 11,027,957.95 & 39,352.86 & 6,461.68 & 5,987.72 \\
问题3 & 上月选择纯历史 & 11,132,033.28 & 37,016.85 & 6,461.68 & 5,987.72 \\
问题4-3 & 固定1月参数 & 11,695,213.43 & 23,070.56 & 6,517.67 & 5,997.11 \\
问题4-3 & 上月选择融合 & 11,711,534.09 & 39,935.96 & 6,517.67 & 5,997.11 \\
问题4-3 & 上月选择纯历史 & 11,812,014.14 & 46,519.39 & 6,517.67 & 5,997.11 \\
\bottomrule
\end{tabular}
\par\smallskip
\begin{minipage}{0.97\linewidth}\footnotesize 三条路径共用4月1日期初库存，此后逐月连续；SOC单位为kWh。本表覆盖275天，不与2—12月334天总费直接相减。\end{minipage}
\end{table}
\par

\FloatBarrier
问题三和4-3的逐月融合方案比固定参数分别多支出2109.14元、16320.66元，9个月中仅3个月、1个月更省；相对逐月历史基准则分别节省104075.33元、100480.05元。同年滚动评价仍属于再分析，不能据月度胜率推出统计显著性或未见年度泛化。
